\section{Discussion} \label{discussion}

We have shown that two agents suffice to optimally explore cycles for the energy model by modifying the \textit{ALE} strategy by Higashikawa et al.\ \cite{Higashikawa2012}
to create the \textit{AMP} strategy, which we also used to achieve optimal online exploration strategies on tadpole graphs with 3 agents for the energy and 4 agents for the time model.
An interesting result for the exploration of tadpole graphs with multiple agents is, that while on cycles online multi-agent strategies with an optimal competitive ratio of $1.5$ 
only achieve a worse competitive ratio than an optimal single-agent exploration strategy,
which reaches a competitive ratio of $\frac{1 + \sqrt{3}}{2} \simeq 1.366$ \cite{MIYAZAKI2009}, better competitive ratios
than in the single-agent case can be achieved on tadpole graphs, with an optimal competitive ratio of $1.5$ with four agents, versus $2$ with a single agent \cite{Brandt2020}. 
In the Appendix we sketch how the exploration strategies on tadpole graphs can be extended onto a graph class we call $n$-tadpole graphs, which are generalized variants of tadpole graphs, where $n$ tails are attached to a cycle instead of only one, with cycles being considered as $0$-tadpole graphs.

\paragraph*{Outlook}
We investigated cycles and tadpole graphs, and sketched extensions of those strategies for $n$-tadpole graphs
in the Appendix. 
These graph classes are all subclasses of unicyclic graphs.
Thus, one possibility for further research could be analyzing how exploration strategies for multi-agent exploration of unicyclic graphs can be achieved. But there are still other graph classes, which could also be considered for multi-agent graph exploration,
like directed graphs or cactus graphs.\\
Moreover, as seen in the results, there is still a gap between the upper and lower bounds in some scenarios.
Hence, another possibility for further research could be proving better lower bounds for tadpole graphs, or finding strategies matching the lower bound with only two or three~agents.\\
Lastly, we considered global communication between all agents, which allowed the identification
of the cycle edge containing the midpoint of the cycle. The strategies presented in this work do
not work if agents cannot communicate their already traversed distance. Thus, an open question is, how the competitive
ratio for the exploration of cycles and tadpole graphs is affected when restricting the communication between the agents.